\raggedright
\chapterpage{arxiv-scope}{STAGING / GELTUNG}{Quellenstand und Claim-Grenzen}{Diese Satzableitung ist ein neuer, nicht eingereichter Review-Kandidat. Sie übernimmt die historische Gesamtausgabe und deren ausdrücklich begrenzte Claim-Disposition.}
\begin{boundary}{Verbindliche Lesart}
Der Manuskriptkörper beschreibt den eingefrorenen Quellenstand vom 3. Oktober 2026. Bezeichnungen wie aktuell und Draft sind historische Aussagen über ihre jeweiligen Quellenzeitpunkte. Keine gegenwärtige Repository-, PR-, Runtime- oder Publikationslage folgt allein aus diesem Satz. Die Original-PDF und das 36-Dateien-ZIP bleiben unverändert; diese neue PDF übernimmt deren Upload-Autorisierung nicht.

Die folgende Matrix berichtet die bestehende Disposition; sie behauptet weder eine neue Replikation historischer Trials noch eine semantische Vollprüfung aller natürlichen Sprachsätze. Kein Claim ist FORMAL\_PROVED. Die normativen Verträge sind kein universeller Vollzugsbeweis. Neue Autorentscheidung, Zielmetadaten und action-time Upload-Freigabe bleiben erforderlich.
\end{boundary}
{\tighttable
\begin{longtable}{@{}L{8mm}L{45mm}L{111mm}@{}}
\toprule ID & Epistemischer Status & Gebundene Aussage \\ \midrule
\endhead
C01 & \code{NORMATIVE} & QIK-VRT wird als Gesamtsystem, nicht als einzelnes Modell beschrieben. \\
C02 & \code{SOURCE_BOUND} & Zwecksetzung, Wiederverwendung und kontinuierliche Verbesserung sind im Entrypoint dokumentiert. \\
C03 & \code{NORMATIVE} & Die Mittel können sich ändern, während ein Vertragskern erhalten werden soll. \\
C04 & \code{INTERPRETATIVE} & Induktiver Invariantenerhalt gilt unter Basis- und Übergangsannahme. \\
C05 & \code{OPEN} & Universelle fehlerfreie Wahrhaftigkeitsdurchsetzung ist nicht nachgewiesen. \\
C06 & \code{SOURCE_BOUND} & EFFECT\_ACK-DONE im I-D ist Freigabeeignung vor dem Effekt, nicht automatisch dessen Vollzug. \\
C07 & \code{INTERPRETATIVE} & Hashbindung ist weder Authentisierung noch Beweis der Quelle. \\
C08 & \code{NORMATIVE} & Frühere Evidenz darf nicht stillschweigend Nachfolger akzeptieren. \\
C09 & \code{INTERPRETATIVE} & Wrapper existiert unter gegebenen berechenbaren Adaptern und Prüfern. \\
C10 & \code{INTERPRETATIVE} & Vollständige Beobachtbarkeit und beliebige Zielerreichbarkeit folgen nicht. \\
C11 & \code{INTERPRETATIVE} & Black-Box-Integration garantiert keine vollständige Identifikation. \\
C12 & \code{NORMATIVE} & Funktionales Bewusstsein ist ein definiertes Zielprofil. \\
C13 & \code{OPEN} & Vollständige Erfüllung / phänomenales Erleben von QIK-VRT. \\
C14 & \code{INTERPRETATIVE} & Zwei Nodes beweisen keine methodische Unabhängigkeit. \\
C15 & \code{EMPIRICALLY_EVIDENCED} & Neun Mikrobenchmark-Mediane sind aus vorhandenen Rohzahlen reproduzierbar. \\
C16 & \code{SOURCE_BOUND} & Historischer Strukturquotient 3.928/20 = 196,4. \\
C17 & \code{SOURCE_BOUND} & Begrenzte 1/2/4-Prozess-Beschleunigung ist dokumentiert. \\
C18 & \code{SOURCE_BOUND} & Der archivierte A/B-Diagnoselauf auf a31ebb7e dokumentiert Credential-Delivery-HOLD. \\
C19 & \code{SOURCE_BOUND} & Der archivierte Messcode startet den Timer erst nach den Repository-Abrufen. \\
C20 & \code{EMPIRICALLY_EVIDENCED} & Textdecodierung/Binärdaten und Modusvergleich sind nicht äquivalent. \\
C21 & \code{OPEN} & Eine vollständige verteilte Beschleunigung ist weiterhin nicht nachgewiesen. \\
C22 & \code{SOURCE_BOUND} & PRs \#437/\#448/\#449 haben neue Kandidaten mit dokumentierten Tests. \\
C23 & \code{SOURCE_BOUND} & Windows-11-AMD64-Abnahme und reale Unterbrechung bleiben offen. \\
C24 & \code{SOURCE_BOUND} & Das Runtime-Manifest auf dem archivierten Main meldet keine gebündelte Windows-Python-Runtime. \\
C25 & \code{INTERPRETATIVE} & Zeitmodell und Kalenderkorridore. \\
C26 & \code{INTERPRETATIVE} & 2,5× lokaler Faktor ist nicht 2,5× Gesamtdauer. \\
C27 & \code{SOURCE_BOUND} & Der v3-Pfad ist technisch implementiert und geprüft, bleibt ohne detached Vertragsaktivierung für Produktion gesperrt. \\
C28 & \code{OPEN} & Zenodo-Publikation und Replikation der Gesamtausgabe auf allen Nodes sind nicht nachgewiesen. \\
C29 & \code{INTERPRETATIVE} & Historische Größe der Entdeckung. \\
C30 & \code{OPEN} & Eine absolute Garantie, Inhalte niemals mehr zu verlieren, ist nicht nachgewiesen. \\
\bottomrule\end{longtable}}
\subsection{Unveränderte methodische Präzisierungen}
\subsection{C18}
Der archivierte A/B-Diagnoselauf auf a31ebb7e dokumentiert Credential-Delivery-HOLD.

Historischer Run 37128318275 / Job 111218106133; kein aktueller Integrationsnachweis und kein Speedup.

Änderungsgrund: Historischen Credential-HOLD an sein ursprüngliches Run-/HEAD-/TREE-Subjekt gebunden.

\subsection{C19}
Der archivierte Messcode startet den Timer erst nach den Repository-Abrufen.

Exakter historischer Blob \code{acafd1ae37dcb9ffcab371fa79e3e874cfe83b91}; lokale Vergleiche, kein verteilter Mutation/Reconciliation/Readback-Zyklus.

Änderungsgrund: Timergrenze, Binärpfad und Git-Modusgrenze auf den konkret geprüften Quellcode begrenzt.

\subsection{C21}
Eine vollständige verteilte Beschleunigung ist weiterhin nicht nachgewiesen.

OPEN: Kein zulässiger positiver Claim aus Mikrobenchmark, Strukturquotient oder erfolgreichem Diagnoseworkflow.

Änderungsgrund: Offenen verteilten Speedup ausdrücklich als nicht nachgewiesen formuliert.

\subsection{C24}
Das Runtime-Manifest auf dem archivierten Main meldet keine gebündelte Windows-Python-Runtime.

Historischer Main \code{8153c4e69a6245aadf67f0741fb50d9707c19e34} / Blob \code{bd2b893176d93814f38d024b5a095b770eee2aa3}; keine Aussage über spätere Runtime-Kandidaten.

Änderungsgrund: Runtime-Zustand als historische Dateibeobachtung statt beweglichem aktuellem Zustand gebunden.

\subsection{C27}
Der v3-Pfad ist technisch implementiert und geprüft, bleibt ohne detached Vertragsaktivierung für Produktion gesperrt.

Integrationsinput 80e34e4d; unveränderte v1/v2-Verträge; REVIEW\_REQUIRED. Technische Readiness erteilt weder Produktionsaktivierung noch AUTHORIZE\_EXACT\_UPLOAD.

Änderungsgrund: Zwischenzeitliche v3-Implementierung von Produktionsaktivierung und späterer Einzelautorisierung getrennt.

\subsection{C28}
Zenodo-Publikation und Replikation der Gesamtausgabe auf allen Nodes sind nicht nachgewiesen.

OPEN: Repository-Materialisierung im Draft-PR ist kein Zenodo-Effekt, keine Main-Promotion und keine all-node acceptance.

Änderungsgrund: Repository-Übernahme von Zenodo-Publikation und Replikationsabnahme getrennt.

\subsection{C30}
Eine absolute Garantie, Inhalte niemals mehr zu verlieren, ist nicht nachgewiesen.

OPEN: Hashbindung und Repository-Persistenz ersetzen keine unabhängigen Replikations- und Restore-Tests.

Änderungsgrund: Absolute Verlustfreiheit als offene Garantie abgegrenzt.

Weitere unveränderte Grenzen: Wahrhaftigkeit ist ein normativer Vertrag und kein universeller Vollzugsbeweis. Bedingte Invarianten- und Wrapper-Argumente sind INTERPRETATIVE; kein Claim ist FORMAL\_PROVED. EFFECT\_ACK vor einem Effekt bleibt von nachträglicher Aufgabenabnahme getrennt. C15 prüft neun vorhandene Zahlenreihen nach, ohne Zeitversuche zu wiederholen. C17 berichtet ausschließlich die dokumentierten 1/2/4-Prozess-Mediane 0,349138 / 0,214989 / 0,163044 Sekunden; die Quotienten 1,623981 und 2,141373 sind Rechenwerte und keine neue Trial-Replikation. C20 reproduziert ausschließlich zwei isolierte Methoden-Kontrollen. C25/C26 sind bedingte Engineering-Szenarien ab tatsächlich verfügbaren Voraussetzungen T0; 7.–11. Oktober ist keine Auslieferungszusage. Subjektives Erleben, globale Systemabnahme, SSO/Biometrie, physische FPGA-Ausführung und methodisch unabhängige Skalierung werden nicht bestätigt.

Die ursprünglichen editorialen Proof-/Return-DRAFTS im ZIP und original-review sind historische Dokumente und dürfen keine Produktion freigeben. Normative Disposition dieser Übernahme: \code{CLAIM_MATRIX.json}, \code{MACHINE_PROOF_BUNDLE.json} und \code{PREPUBLICATION_RETURN_RECEIPT.json}. Die Produktionsaktivierung und die spätere einmalige AUTHORIZE\_EXACT\_UPLOAD-Entscheidung sind detached und bleiben abwesend.

\subsection{Zwischenzeitlicher Methoden-Nachfolger}
Während dieser Übernahme wurde \code{a6f930bb9c8bad58a10d17cd10c87d94c38f0b25} / TREE \code{fbfdba90e086128b5acd4ce5602f0abefed8fce4} als separater Methodenreparatur-Kandidat materialisiert. Dessen Messcode (Blob \code{0d8e10d419e453a832453343b8a82b330193c330}) verwendet Rohbytes und NUL-Pfade, prüft Git-Modus/Typ/Objekt/Inhalt und trennt lokale Vergleichskosten vom vollständigen Remote-Fetch/isolierte-Replik-Abgleich/frischen-Readback-Zyklus. Die Reconciliation verändert ausschließlich die lokale Wegwerf-Replik. Sie ist keine Mutation der beiden Produktionsremotes und beweist keinen verteilten Produktions-Speedup. Die historischen C18/C19-Befunde in der Original-PDF werden nicht umgeschrieben; die neue technische Prüfung muss den integrierten Nachfolger ausführen.

